\documentclass[11pt]{article}
\usepackage[a4paper,margin=19mm]{geometry}
\usepackage{fontspec}
\setmainfont{lmroman10-regular.otf}
\usepackage{amsmath,amssymb,amsthm,mathtools,xcolor,graphicx,hyperref}
\hypersetup{hidelinks}
\numberwithin{equation}{section}
\newtheorem{theorem}{Theorem}[section]
\newtheorem{lemma}[theorem]{Lemma}
\newtheorem{proposition}[theorem]{Proposition}
\newtheorem{corollary}[theorem]{Corollary}
\newtheorem{claim}[theorem]{Claim}
\newtheorem{fact}[theorem]{Fact}
\newtheorem{conjecture}[theorem]{Conjecture}
\theoremstyle{definition}
\newtheorem{definition}[theorem]{Definition}
\newtheorem{example}[theorem]{Example}
\newtheorem{remark}[theorem]{Remark}
\newtheorem{question}[theorem]{Question}
\theoremstyle{plain}
\newcommand{\C}{\mathbb C}
\newcommand{\R}{\mathbb R}
\newcommand{\Q}{\mathbb Q}
\newcommand{\Z}{\mathbb Z}
\newcommand{\N}{\mathbb N}
\newcommand{\eee}{\mathrm e}
\newcommand{\cclass}[1]{\textsf{#1}}
\newcommand{\cproblem}[1]{\textsc{#1}}
\newcommand{\ie}{i.\,e.}
\newcommand{\Ie}{I.\,e.}
\newcommand{\eg}{e.\,g.}
\newcommand{\Eg}{E.\,g.}
\newcommand{\phd}{Ph.\,D.}
\newcommand{\msc}{M.\,S.}
\newcommand{\bsc}{B.\,S.}
\newcommand{\expref}[2]{\hyperref[#2]{#1~\ref{#2}}}
\newcommand{\expeqref}[2]{\hyperref[#2]{#1~\eqref{#2}}}
\newcommand{\exprefsilent}[2]{\hyperref[#2]{#1}}
\providecommand{\etal}{et~al.}
\usepackage{eprint}
\usepackage{amsmath,mathtools, amsthm, amssymb}
\usepackage{cleveref,enumerate}
%


\newcommand{\E}{\mathbb{E}}
\newcommand{\F}{\mathbb{F}}
\newcommand{\FF}{\mathcal{F}}
\newcommand{\LL}{\mathcal{L}}
\newcommand{\BB}{\mathcal{B}}
%
\DeclareMathOperator{\perm}{perm}           %
\newcommand{\BBperm}{\mathcal{B}_{\perm}}    %
\newcommand{\eps}{\varepsilon}
%
\DeclareMathOperator{\sign}{sign}     %
\newcommand{\Lone}[1]{\textbf{L1}(#1)}
%
\DeclareMathOperator{\poly}{poly}      %
\newcommand{\zero}{\overline{0}}

\newtheorem{open}[theorem]{Open problem}
\title{Fourier-level L1 bounds for low-degree binary phase polynomials}
\author{Eshan Chattopadhyay, Pooya Hatami, Kaave Hosseini, Shachar Lovett}
\date{doi:10.4086/toc.2019.v015a010; CC BY 3.0}
\begin{document}
\maketitle
\noindent\textit{Editorial note.} Complete original author body and bibliography retained. Restriction-closure, symmetry/noticability and low-degree binary polynomial conventions preserved. Explicit author limitation concerning low-degree-polynomial PRGs retained.\par\medskip
\noindent\textit{Allowed standard prerequisites (editorial).}\begin{itemize}
\item Boolean Fourier analysis, random restrictions and elementary conditional probability. Original definitions and all local polarization, rounding or Fourier recurrence arguments retained.. No collector reproof; applications/open problems remain their original scope.
\end{itemize}\medskip
\section{Introduction}

Pseudorandom generators (PRG) are widely studied in complexity theory. There are several general frameworks used to construct PRGs.
One is based on basic building blocks, such as
%
small-bias   %
generators~\cite{naor1993small,alon1992simple}, $k$-wise independence,
or expander graphs~\cite{hoory2006expander}. Another approach is based on
the %
hardness vs.\ randomness paradigm, which was introduced by
Nisan and Wigderson~\cite{nisan1988hardness} and has been very influential. Many of the hardness results used in the latter framework
are based on random restrictions, and the analysis of how
they simplify the target class of functions.
%
The number of papers in these lines of work is on the
order of hundreds, so we do not even attempt to give a comprehensive survey of them all; instead we refer the reader to survey articles \cite{goldreich2010primer,vadhan2012pseudorandomness}.

The purpose of this paper is to introduce a new framework for constructing PRGs based on polarizing random walks. We develop the theory
in this paper and give a number of applications; perhaps the most notable one is a PRG for functions of sensitivity $s$ whose seed
length is polynomial in $s$. But, as this is a new framework, there are many questions that arise, both technical and conceptual,
and we view this paper as mostly preliminary, with the hope that many more applications would follow.

\subsection{PRGs and fractional PRGs}
Let $\FF$ be a class of Boolean functions $f:\{-1,1\}^n \to \{-1,1\}$. The standard definition of a PRG for $\FF$ with error $\eps>0$,
is a random variable $X \in \{-1,1\}^n$ such that
$$
\forall f\in \FF, |\E_X[f(X)] - \E_U[f(U)]| \le \eps,
$$
where $U$ denotes a random variable with the uniform distribution in $\{-1,1\}^n$.  We relax this definition by introducing a new object called a \emph{fractional PRG}, defined in the next paragraph.

To prepare the notation for the definition, identify $f$ with a real multilinear polynomial, namely its Fourier expansion. This extends $f:\{-1,1\}^n\rightarrow \{-1,1\}$ to $f:\R^n \to \R$, although, we would only
be interested in inputs from $[-1,1]^n$. Observe that if $x \in [-1,1]^n$ then $f(x) = \E_X[f(X)]$ where $X \in \{-1,1\}^n$ is a random variable
sampled as follows: for every $i \in [n]$ sample $X_i \in \{-1,1\}$ independently with $\E[X_i]=x_i$.
In particular, $f$ on $[-1,1]^n$ is bounded, namely $f:[-1,1]^n \to [-1,1]$. Also, $f(\zero) = \E_U[f(U)]$. The following is a key definition.

\begin{definition}[Fractional PRG]
Let $f:[-1,1]^n \to [-1,1]$ be real multilinear. A fractional PRG for $f$ is a random variable $X \in [-1,1]^n$ such that
$$
|\E_X[f(X)] - f(\zero)| \le \eps.
$$
	Moreover, $X$ has seed length $r$ if $X = G(U)$ for some function $G:\{-1,1\}^r \to [-1,1]^n$.
\end{definition}

One trivial construction of a fractional PRG is $X\equiv \zero$ but this is not going to be useful for our purpose of constructing PRGs.
To disallow such examples, we require each coordinate of $X$ to be far from zero with some noticeable probability. Formally, $X \in [-1,1]^n$ is called $p$-noticeable
if $\E[X_i^2] \ge p$ for all $i=1,\ldots,n$.

A good example to keep in mind is the following. Let $G: \{ -1,1\}^r \rightarrow \{-1,1\}^{n}$ be a (Boolean valued) function,
and set $X=p G(U)$, where $U \in \{-1,1\}^r$ is uniform. Notice that $X$ is $p^2$-noticeable.
%
%

Fractional PRGs are easier to construct than standard PRGs, as they can take values in $[-1,1]^n$. For example, assume that $f$ has Fourier tails
bounded in $L_1$. That is,  there exist parameters $a,b \ge 1$ for which
$$
\sum_{S \subseteq [n]: |S|=k} |\hat{f}(S)| \le a \cdot b^k \qquad \forall k=1,\ldots,n.
$$
We show (in \expref{Lemma}{lemma:L1}) that if $X \in \{-1,1\}^n$ is roughly $(\eps/a)$-biased, then $pX$ is a fractional PRG for $f$ with $p \approx 1/b$ and error $\eps$. The reason is that
this choice of $p$ controls all the Fourier coefficients of $f$ with large Hamming weight, while $X$ controls the ones with small weight.
(In fact, to optimize parameters one can choose $X$ to be almost $k$-wise independent; see \expref{Lemma}{lemma:L1} for details).
In any case, note that $pX$ is $p^2$-noticeable as $pX$ takes values in $\{-p,p\}^n$.

\subsection{Fractional PRG as steps in a random walk}

Let $X \in [-1,1]^n$ be a fractional PRG for $f$ with error $\eps$. That is,
$$
|\E_X[f(X)] - f(\zero)| \le \eps.
$$


The goal is to construct a random variable $Y \in \{-1,1\}^n$ such that $\E_Y[f(Y)] \approx f(\zero)$, where the fractional PRG $X$ provides
a ``small step'' towards this approximation. If we can combine these small steps in a way that they converge fast to $\{-1,1\}^n$, then we would be done. To be a bit more precise, consider a random walk starting at $\zero$  with the following properties:
\begin{enumerate}
\item The value of $f$ at each step on average does not change by too much.
\item The random walk converges fast to $\{-1,1\}^n$.
\end{enumerate}

Observe that if we take $X$ as the first step, then property 1 is satisfied for the first step. Considering later steps leads to the following question: Given a point $y\in [-1,1]^n$, can we find a random variable $A=A(y,X)$ such that
$$
|\E[f(A)] - f(y)| \le \eps,
$$
and such that $A$ takes values closer to Boolean values?
We show that this is indeed the case if we assume that $X$ not only fools $f$, but also fools any possible restriction of $f$.

To formalize this, let $\FF$ be a family of $n$-variate Boolean functions $f:\{-1,1\}^n \to \{-1,1\}$. We say that $\FF$ is closed under restrictions if for any
$f \in \FF$, if we fix some inputs of $f$ to constants $\{-1,1\}$, then the new restricted function is still in $\FF$. Most natural families of Boolean functions studied
satisfy this condition. Some examples are functions computed by small-depth circuits, functions computed by bounded width branching programs, and functions
of low sensitivity.

We show that if $X$ is a fractional PRG for such $\FF$, then it can be used to approximate $f(y)$ for any $y \in [-1,1]^n$. Define $\delta_y \in [0,1]^n$ by
$(\delta_y)_i = 1 - |y_i|$. For $x,x' \in [-1,1]^n$ define $x \circ x' \in [-1,1]^n$ to be their coordinate-wise product, $(x \circ x')_i = x_i x'_i$.
Note that under this definition, the subcube $\{y + \delta_y \circ x: x \in [-1,1]^n\}$ is the largest symmetric subcube of $[-1,1]^n$ centered at $y$.

We show (\expref{Claim}{claim:prg_around_nonzero}) that if $X \in [-1,1]^n$ is a fractional PRG for $\FF$ which is closed under restrictions, then for any
$f \in \FF$ and any $y \in [-1,1]^n$ it holds that
$$
|\E[f(y + \delta_y \circ X)] - f(y)| \le \eps.
$$
Technically, we need to also assume that $X$ is \emph{symmetric}, which means that $\Pr[X=x]=\Pr[X=-x]$ for all $x$. This is easy to achieve
from any $X$ which is not symmetric, for example by multiplying
$X$ with a uniform bit (thus, increasing its seed length by $1$ bit).

\subsection{Polarization and fast convergence}
Our next goal is to show fast convergence of the random walk to $\{-1,1\}^n$. To that end, we need
to analyze the following martingale:
\begin{align*}
&Y_1=X_1,\\
&Y_i=Y_{i-1} + \delta_{Y_{i-1}} \circ X_i,
\end{align*}
where $X_1,X_2,\ldots$ are independent copies of a fractional PRG\@. We show that for some $t$ not too large, $Y_t$ is close to a point in $\{-1,1\}^n$.
But why would that be true? This turns out to be the result of \emph{polarization} in the random walk. It suffices to show this for every coordinate individually.

So, let $Z_1,Z_2,\ldots \in [-1,1]$ be independent random variables (which are the $i$-th coordinate of $X_1,X_2,\ldots$ for some fixed $i$), and define the following one-dimensional martingale:
\begin{align*}
&W_1=Z_1,\\
&W_i=W_{i-1} + (1-|W_{i-1}|) Z_i.
\end{align*}
\expref{Claim}{claim:polar_onedim} shows that if (i) $Z_i$ is symmetric, and (ii) $\E[Z_i^2] \ge p$ (which follows from our assumption that the fractional PRG is $p$-noticeable), then it holds that
$$
\Pr[|W_t| \ge 1-\delta] \ge 1-\delta
$$
for $t=O(\log (1/\delta) / p)$. Setting $\delta=\eps/n$ guarantees that with probability $1-\eps$ all the coordinates of $Y_t$ are $\eps/n$ close to $\{-1,1\}$.
Then a simple argument shows that rounding the coordinates gives a PRG with error $O(\eps)$, as desired.

We now state our main theorem.

\begin{theorem}[Main theorem, informal version of \expref{Theorem}{thm:main}]
\label{thm:main_informal}
Let $\FF$ be a family of $n$-variate Boolean functions that is closed under restrictions.
Let $X \in [-1,1]^n$ be a symmetric $p$-noticeable fractional PRG for $\FF$ with error $\eps$.
Set $t=O(\log(n/\eps)/p)$ and let $X_1,\ldots,X_t$ be i.i.d.\ copies of $X$.
Define the following random variables taking values in $[-1,1]^n$:
$$
Y_0=\zero; \qquad Y_i = Y_{i-1} + \delta_{Y_{i-1}} \circ X_i \qquad i=1,\ldots,t.
$$
Let $G=\sign(Y_t) \in \{-1,1\}^n$ obtained by taking the sign of the coordinates in $Y_t$.
Then $G$ is a PRG for $\FF$ with error $(t+1)\eps$.
\end{theorem}

Note that computing this PRG only involves basic operations such as addition and multiplication over the reals with bounded error. 
\subsection{PRG for functions with bounded Fourier tails}
As mentioned above, the families of Boolean functions that are fooled by our PRG include ones that satisfy the following two properties: (i) being closed under restrictions; (ii) having bounded $L_1$
Fourier tails. Tal~\cite{tal2017tight} showed that the latter condition follows from a widely studied condition, that of bounded $L_2$ Fourier tails.
Thus, using existing bounds for $L_2$ Fourier tails, we get that our PRG fools several classes of Boolean functions.
Below we list the results for error $\eps=O(1)$, and refer the reader to the corresponding claims for the details of the full range of parameters:

\begin{enumerate}
%
%


\item {\bf Functions of sensitivity $\mathbf{s}$}: seed length $O(\mathbf{s}^3 \log \log n)$. The best previous construction due to Hatami and Tal~\cite{hatami2017pseudorandom} required seed length subexponential in $\mathbf{s}$ (concretely, their dependence on
  %
  %
  %
  %
   $\mathbf{s}$
  %
  is $\exp(\sqrt{\mathbf{s}})$). See \expref{Corollary}{cor:sensitivity} for details.

%
%

\item {\bf Unordered read-once branching programs of width $\mathbf{w}$}: seed length $O(\log^{2\mathbf{w}+1} n \cdot \log\log n)$. This is quadratically worse than the best known PRG due to Chattopadhyay et al.~\cite{chattopadhyay2018improved}. However, our PRG construction does not utilize the branching program structure at all, except to obtain the Fourier tail bounds. See \expref{Corollary}{cor:ROBP} for details.

%
%

\item {\bf Permutation unordered read-once branching programs of width $\mathbf{w}$}: seed length $O(\mathbf{w}^4 \log n  \cdot \log\log n)$. This improves the dependence on $n$ quadratically compared to the previous best PRG due to Reingold et al.~\cite{reingold2013pseudorandomness}. See \expref{Corollary}{cor:perm_BP} for details.
	
%
\item {\bf Bounded-depth circuits}:  %
if $f$ is computed by $\cclass{AC}^0$ circuits of depth $\mathbf{d}$ and size $\poly(n)$, our PRG has seed length
%
%
	$O(\log^{2\mathbf{d}-1} n \cdot \log\log n )$. This is quadratically worse than the best known PRG due to Tal~\cite{tal2017tight}.
See \expref{Corollary}{cor:AC0} for details.
\end{enumerate}

Other than the PRG for functions of low sensitivity, all the other PRGs are comparable to the best known tailored PRG\@. However, the main message
is that {\bf they are all the same PRG}. Our general theorem is the following.

\begin{theorem}[PRG for functions of bounded $L_1$ Fourier tail, informal version of \expref{Theorem}{thm:L1}]
\label{thm:L1_informal}
Let $\FF$ be a family of $n$-variate Boolean functions closed under restrictions.
Assume that there exist $a,b \ge 1$ such that for every $f \in \FF$,
$$
\sum_{S \subseteq [n]: |S|=k} |\hat{f}(S)| \le a \cdot b^k.
$$
Then, for any $ \eps\leq {1}/{\poly(b \log n)}$ there exists an explicit PRG $X \in \{-1,1\}^n$ which fools $\FF$ with error $\eps>0$,
whose seed length is $O(\log(n/\eps)  (\log\log n + \log(a/\eps)) b^2)$.
\end{theorem}

%
%
We note again that by a result of Tal~\cite{tal2017tight}, \expref{Theorem}{thm:L1_informal} holds also if we instead assume a bound on the $L_2$ Fourier tails
(which are more common), namely if we assume that for every
$f \in \FF$ it holds that
$$
\sum_{S \subseteq [n]: |S| \ge k} \hat{f}(S)^2 \le a \cdot 2^{-k/b}.
$$

\subsection{PRG for functions which simplify under random restriction}
A major component in prior constructions of PRGs that are based on random restrictions is finding a much smaller set of ``pseudorandom retrictions.'' Ajtai and Wigderson \cite{ajtai1985deterministic} proposed such a PRG for
%
low-depth circuits.   %
%
Much follow-up work has been  %
based on this framework to build PRGs for various classes of functions including
%
low-depth      %
circuits, branching programs, low-sensitivity functions \cite{trevisan2013derandomized,gopalan2012better,reingold2013pseudorandomness,chattopadhyay2018improved, hatami2017pseudorandom}, and  a major component of the analysis is proving that the derandomized random restrictions work.

Our framework for constructing PRGs directly applies to function families that simplify under random restrictions without the need to derandomize the restrictions.  Let $\FF$ be a family of functions  $f:\{-1,1\}^n \to \{-1,1\}$  which are extended multilinearly to $[-1,1]^n$. Fix a parameter $0<p<1$ and define the $p$-averaged function of $f$, denoted $f_p:\{-1,1\}^n \to [-1,1]$, as follows:
sample $A \subseteq [n]$ where $\Pr[i \in A]=p$ independently for $i \in [n]$, and define
$$f_p(x) = \E_{A,U} [f(x_A,U_{A^c})]$$
where $x_A \in \{-1,1\}^A$ is the restriction of the input $x$ to the coordinates in $A$, and $U \in \{-1,1\}^n$ is independently and uniformly chosen. The crucial observation (\expref{Claim}{claim:p_averaged})
is that for every $x \in \{-1,1\}^n$ it holds that
$$
f(px) = f_p(x).
$$

Suppose now we have a standard PRG $X$ for the class of $p$-averaged functions $\FF_p = \{f_p : f\in \FF\}$. Note a PRG for the $p$-random restriction of functions in $\FF$ would do,
%
%
as $f_p$ is a convex combination of $p$-random restrictions of $f$ (namely, averaging over $U$). Then, using our observation above,  this implies that $X'=pX$ is a fractional PRG for the class $\FF$. Now by using our framework of viewing this fractional PRG as a random walk step, one can derive a standard PRG for $\FF$ using $O(\log(n/\epsilon)/p^2)$ independent copies of $X$.

\subsection{Fourier tails of low-degree \texorpdfstring{$\F_2$}{F2} polynomials}
Viola~\cite{viola2009sum} gave a construction of a pseudorandom generator which fools $n$-variate polynomials over $\F_2$.
The construction is the XOR of $d$ independent small-bias generators.
We wonder whether our framework can be used to achieve
similar bounds. In particular, we raise the following problem: does the class of low-degree polynomials over $\F_2$ have bounded
$L_1$ Fourier tails? It's trivially true for $d=1$ and it can be shown to hold for $d=2$. However, to the best of our knowledge nothing was known for
 $d \ge 3$.

We show (see \expref{Theorem}{thm:polynomialspectralnorm}  for more details) that for any Boolean function $f: \{ -1,1\}^{n} \rightarrow \{-1,1 \}$ computed by a $\F_2$-polynomial of degree at most $d$, the following $L_1$ Fourier tail bound holds: $$ \sum_{|S|=k} |\widehat{f}(S)| \le k^k 2^{3dk} \qquad \forall k=1,\ldots,n.$$

This bound however falls short of  implying a PRG using our techniques, and we conjecture that the correct bound is $c_{d}^{k}$, for some constant $c_d=2^{O(d)}$.

\subsection{PRGs with respect to arbitrary product distributions} We note the following interesting generalization of our results that  is almost direct from our techniques. Consider the problem of ``fooling'' a family of functions with respect to an arbitrary product distribution $D$ on $\{ -1,1\}^n$ (the uniform distribution being a special case). More formally, given a distribution $D$ on $\{-1,1 \}^n$ and a family of functions $\FF$, we say that a random variable $X$ is a PRG for $\FF$ (with respect to $D$) if $|\E[f(D)]  - \E[f(X)]| \le \epsilon$.

We show a way to fool functions with respect to arbitrary product distributions.
\begin{corollary}
\label{cor:arbitrary_product_fool}
Let $\FF$ be a family of $n$-variate Boolean functions which is closed under restrictions and let $D$ be any product distribution on $\{-1,1 \}^n$.
Let $X \in [-1,1]^n$ be a symmetric $p$-noticeable fractional PRG for $\FF$ with error $\eps$ and seed length $\ell$. Let $t= O(\log(n/\epsilon)/p)$. Then there exists an explicit  PRG for $\FF$ with respect to $D$ with error $t\eps$ and seed length $t \ell$.
\end{corollary}
\begin{proof}[Proof sketch]If $D$ is a product distribution on $\{-1,1\}^n$, then $\E[f(D)] = f(\alpha)$, where $\alpha=\E[D] \in [-1,1]^{n}$. Thus, we now start our random walk (defined by the fractional PRG) from the point $\alpha$ instead of from $\zero$, and the convergence follows from polarization in exactly the same way.
\end{proof}
Thus all our PRG results in fact generalize to PRGs with respect to arbitrary product distributions. To the best of our knowledge, we are not aware of any non-trivial PRGs against arbitrary product distributions  for the classes of functions we study. We wonder if this notion of fooling arbitrary product distributions has interesting applications.

%
\subsection{Related work}
%
%
The line of research closest in spirit to our paper, and which motivated our result, is that of using random and pseudo-random restrictions
%
%
to construct PRGs. A good example is due to Gopalan et al.~\cite{gopalan2012better} which uses pseudo-random restrictions to construct PRGs. Our framework
can be seen as extending this, as we do not need to analyze pseudo-random restrictions; instead, we analyze fractional PRGs, where the restriction
happens automatically from the fractional PRG structure, and no derandomization is necessary.

Another line of work is the use of random walks in combinatorial optimization, for example in the algorithmic versions of Spencer's theorem~\cite{bansal2010constructive,lovett2012constructive} and
%
follow-up work.     %
It would be interesting to see if polarization can be used
to speed up random walks in combinatorial optimization as well.


\subsection{Open problems}
As we give a new framework for constructing PRGs,
there are many open problems that arise, both conceptual and technical.

\subsubsection{Early termination}
Our analysis requires a random walk with $t=O(\log(n/\eps)/p)$ steps, each coming from a $p$-noticeable fractional PRG\@.
We believe that for some natural families of functions shorter random walks might also suffice, but we do not
know how to show this. We discuss this further in \expref{Section}{section:smoothness}.


\begin{open}
\label{open:early_termination}
Find conditions on classes of Boolean functions so that short random walks can be used to construct PRGs.
In particular, are there nontrivial classes where the number of steps is independent of $n$?
\end{open}

\subsubsection{Less independence}
Our analysis of \expref{Theorem}{thm:main} currently requires to assume $t$ independent copies of a fractional PRG $X$. It might be possible that these copies can be chosen in a less independent form, where the analysis still holds.

\begin{open}
Can the fractional PRGs $X_1,\ldots,X_t$ in \expref{Theorem}{thm:main} be chosen not independently, such that the conclusion still holds? Concrete examples to consider are $k$-wise independence for $k \ll t$, or using an expander random walk.
\end{open}


\subsubsection{More applications}
Our current applications follow from the construction of a fractional PRG for functions with bounded Fourier tails.
The fractional PRG itself follows from standard constructions in pseudo-randomness (almost $k$-wise independent)
adapted to our scenario. It will be interesting to try and find other classes of Boolean functions for which different constructions of fractional PRG work.

\subsubsection{Gadgets}
We can view the random walk as a ``gadget construction.'' Given independent $p$-noticeable fractional PRGs
$X_1,\ldots,X_t \in [-1,1]^n$, view them as the rows of a $t \times n$ matrix, and then apply a gadget $g:[-1,1]^t \to \{-1,1\}$
to each column to obtain the outcome in $\{-1,1\}^n$. We show that the random walk gives such a gadget which converges for $t=O(\log(n/\eps)/p)$.
Many constructions of PRGs can be viewed in this framework, where typically $X_i \in \{-1,1\}^n$. Ours is the first construction which allows
$X_i$ to take non-Boolean values. It is interesting whether other gadgets can be used instead of the random walk gadget, and whether there are
general properties of gadgets that would suffice.

\subsubsection{Low-degree polynomials}
As discussed above, we wonder if our techniques can be used to construct a PRG for low-degree $\F_2$ polynomials. In particular, we ask if one could improve the bounds we obtain (see \expref{Theorem}{thm:polynomialspectralnorm}) on the $L_1$ Fourier tails of low-degree $\F_2$ polynomials.

\begin{open}
 Let $f=(-1)^p$ where $p:\F_2^n \to \F_2$ is a polynomial of degree $d$. Is there a constant $c_d$ such that  $\sum_{S : |S|=k} |\hat{f}(S)| \le c_d^{k}$
which is independent of $n$? In particular, we conjecture
that $c_d=2^{O(d)}$ should work.
\end{open}
Note that the exponential dependence on $k$ is needed,
as witnessed from the following example: consider the quadratic $\F_2$ polynomial
\[
  q(x)=\sum_{i=1}^{n/2}x_{2i-1}x_{2i}.
\]
Then $(-1)^{q}$ has Fourier $L_1$ weight
$\binom{n}{n/2} \cdot 2^{-n/2} = 2^{\Omega(n)}$ on the $(n/2)$-th
level.

%
%
%
%
%
%
%

%
%
%
%
%

%
%


\subsection{Paper organization}
We describe the general framework in detail in \expref{Section}{sec:framework}. We prove \expref{Theorem}{thm:main} in \expref{Section}{sec:proof_main}. We describe
applications in \expref{Section}{sec:spectral}. Our framework also applies to function families that simplify under random restrictions. We describe this in \expref{Section}{sec:random_restriction}. We prove $L_1$ Fourier tail bounds for low-degree $\F_2$ polynomials in \expref{Section}{section:lowdegree_spectral}. We try to partially answer the question related to early termination of the random walk in \expref{Section}{section:smoothness}.

\section{General framework}
\label{sec:framework}

\subsection{Boolean functions}
Let $f:\{-1,1\}^n \to [-1,1]$ be an $n$-variate Boolean function, identified with its multilinear extension,
also known as its Fourier expansion. For $x \in [-1,1]^n$ define $f(x) = \sum_{S \subseteq [n]} \hat{f}(S)x^S$ where  $x^S = \prod_{i \in S} x_i$.
As $f$ is multilinear, a convenient viewpoint is to view $f(x)$ as computing the expected value of $f$ on a product distribution on $\{-1,1\}^n$. That is, let $W=W(x) \in \{-1,1\}^n$ be a random variable, where $W_1,\ldots,W_n$ are independently chosen so that $\E[W_i]=x_i$. Then $f(x)=\E f(W)$. In particular, $f(\zero) = \E f(U)$, where $U \in \{-1,1\}^n$ is uniformly chosen.

A family $\FF$ of $n$-variate Boolean functions
is said to be \emph{closed under restrictions} if for any $f \in \FF$ and any function
$f':\{-1,1\}^n \to \{-1,1\}$ obtained from $f$ by fixing some of its inputs to $\{-1,1\}$ it holds that also $f' \in \FF$.

\subsection{Pseudorandom generators}

Let $\FF$ be a family of $n$-variate Boolean functions. The following is the standard definition of a
pseudorandom generator (PRG) for $\FF$, adapted to our notation.

\begin{definition}[PRG]
A random variable $X \in \{-1,1\}^n$ is a PRG for $\FF$ with error $\eps$, if for any $f \in \FF$ it holds that
$\left| f(\zero) - \E f(X) \right| \le \eps$.
\end{definition}

We introduce the notion of a \emph{fractional PRG}. It is the same as a PRG, except that the random variable is allowed to take  values in $[-1,1]^n$,
instead of only Boolean values. We assume that $X$ has finite support.
 
\begin{definition}[Fractional PRG]
A random variable $X \in [-1,1]^n$ with finite support, is a fractional PRG for $\FF$ with error $\eps$, if for any $f \in \FF$ it holds that
$\left| f(\zero) - \E f(X) \right| \le \eps$.
\end{definition}

Our main goal will be to ``amplify'' fractional PRGs for $\FF$ in order to obtain PRGs for $\FF$. To that end, we need to enforce some non-triviality
conditions on the fractional PRG\@. For example, $X=\zero$ is a fractional PRG for any function. We require that for any coordinate $i \in [n]$,
the value of $X_i$ is far from zero with noticeable probability. Formally, we require a noticeable second moment.

\begin{definition}[$p$-noticeable random variable]
A random variable $X \in [-1,1]^n$ is $p$-noticeable if for every $i \in [n]$, $\E[X_i^2] \ge p$.
\end{definition}

For technical reasons, we would also need $X$ to be \emph{symmetric}, which means that the distribution of $-X$ is the same
as the distribution of $X$. This is easy to achieve, for example by multiplying all elements of $X$ with a uniformly chosen sign.

\subsection{Polarizing random walks}
The main idea is to view a fractional PRG as steps in a random walk in $[-1,1]^n$ that converges to $\{-1,1\}^n$.
To that end, we define a gadget that implements the random walk; and moreover, that allows for fast convergence. As we will see later, the fast convergence is an effect of polarization.


\begin{definition}[Random walk gadget]
For any $t \ge 1$ define the random walk gadget $g_t:[-1,1]^t \to [-1,1]$ as follows.
Let $a_1,\ldots,a_t \in [-1,1]$.
Define $g_1(a_1):=a_1$ and for $t>1$,
$$
g_t(a_1,\ldots,a_t) := g_{t-1}(a_1,\ldots,a_{t-1}) + (1-|g_{t-1}(a_1,\ldots,a_{t-1})|) a_t.
$$
We extend the definition to act on bit-vectors.
Define $g_t^n:([-1,1]^n)^t \to [-1,1]^n$ as follows. For $x_1,\ldots,x_t \in [-1,1]^n$ define
$$
g_t^n(x_1,\ldots,x_t) = \left(g_t(x_{1,1},\ldots,x_{t,1}),\ldots,g_t(x_{1,n},\ldots,x_{t,n})\right).
$$
Equivalently, we can view $g_t^n$ as follows: construct a $t \times n$ matrix whose rows are $x_1,\ldots,x_t$; and
then apply $g_t$ to each column of the matrix to obtain a resulting vector in $[-1,1]^n$.
\end{definition}

The following theorem shows how to ``amplify'' fractional PRGs using the random walk gadget to obtain a PRG\@.
Below, for $x \in [-1,1]^n$ we denote by $\sign(x) \in \{-1,1\}^n$ the Boolean vector obtained by taking
the sign of each coordinate (the sign of $0$ can be chosen arbitrarily).

\begin{theorem}[Amplification Theorem]
\label{thm:main}
Let $\FF$ be a family of $n$-variate Boolean functions which is closed under restrictions.
Let $X \in [-1,1]^n$ be a symmetric $p$-noticeable fractional PRG for $\FF$ with error $\eps$. Set $t=O(\log(n/\eps)/p)$ and let $X_1,\ldots,X_t$
be iid copies of $X$. Define a random variable $G \in \{-1,1\}^n$ as follows:
$$
G:=G(X_1,\ldots,X_t) = \sign(g_t^n(X_1,\ldots,X_t)).
$$
Then $G$ is a PRG for $\FF$ with error $(t+1)\eps$.
\end{theorem}


\section{Proof of Amplification Theorem}
\label{sec:proof_main}
We prove \expref{Theorem}{thm:main} in this section. {}From here onwards, we fix a family $\FF$ of $n$-variate Boolean functions which is closed under restrictions.
The proof is based on the following two lemmas.
The first lemma amplifies a $p$-noticeable fractional PRG to a $(1-q)$-noticeable fractional PRG\@. The second
lemma shows that setting $q=\eps/n$, the latter fractional PRG can be rounded to a Boolean-valued PRG without incurring
too much error.

\begin{lemma}[Amplification lemma]
\label{lemma:amplify}
Let $X_1,\ldots,X_t \in [-1,1]^n$ be independent symmetric $p$-noticeable fractional PRGs for $\FF$ with error $\eps$.
Define a random variable $Y \in [-1,1]^n$ as
$$
Y:=g_t^n(X_1,\ldots,X_t).
$$
Then $Y$ is a $(1-q)$-noticeable fractional PRG for $\FF$ with error $t\eps$, where $q=2^{-\Omega(pt)}$.
\end{lemma}

\begin{lemma}[Rounding lemma]
\label{lemma:round}
Let $Y \in [-1,1]^n$ be a $(1-q)$-noticeable fractional PRG for $\FF$ with error $\eps$.
Then $\sign(Y) \in \{-1,1\}^n$ is a PRG for $\FF$ with error $\eps+qn$.
\end{lemma}

\expref{Theorem}{thm:main} follows directly by applying \expref{Lemma}{lemma:amplify} with $t=O(\log(n/\eps)/p)$ to obtain
$q=\eps/n$ and then applying \expref{Lemma}{lemma:round}.

\subsection{Proof of \expref{Lemma}{lemma:amplify}}
We prove \expref{Lemma}{lemma:amplify} in this section. We need to prove two claims: that $g_t^n(X_1,\ldots,X_t)$ is a fractional
PRG for $\FF$ with error $\eps t$, and that it is $(1-q)$-noticeable. This is achieved in the following sequence of claims.

First we need some notations.
For $y \in [-1,1]^n$ define $\delta_y \in [-1,1]^n$ by $(\delta_y)_i := 1-|y_i|$. For two vectors $x,y \in [-1,1]^n$ define $x \circ y \in [-1,1]^n$ to be their pointwise product, namely $(x \circ y)_i := x_i y_i$.
Observe that $\{y + \delta_y \circ x: x \in [-1,1]^n\}$ is the largest symmetric subcube in $[-1,1]^n$ centered at $y$.

\begin{claim}
\label{claim:prg_around_nonzero}
Let $X \in [-1,1]^n$ be a fractional PRG for $\FF$ with error $\eps$. Then for any $f \in \FF$ and any $y \in [-1,1]^n$,
$$
\left| f(y) - \E f(y + \delta_y \circ X) \right| \le \eps.
$$
\end{claim}

\begin{proof}
Consider a distribution over $F \in \FF$ obtained from $f$ by fixing the $i$-th input to $\sign(y_i)$ with probability $|y_i|$,
independently for each $i$. That is,
$$
F(x) := f(R(x)),
$$
	where $R(x) \in [-1,1]^n$ is a random variable obtained by sampling $R_1,\ldots,R_n$ independently where
	each $R_i$ is chosen as follows. Pick $R_i(x)=\sign(y_i)$ with probability $|y_i|$ and with probability $1-|y_i|$ do as follows: pick $R_i(x) = 1$ with probability $(x_i+1)/2$ and pick $R_i(x)=-1$ otherwise. It's easy to check that $\E_R(R(x)) = y + \delta_y \circ x$. By  multilinearity of $f$, and as $R(x)$ is a product distribution, for all $x\in [-1,1]^n$,
$$
\E_F[F(x)] = \E_R [f (R(x))] = f( \E_R[R(x)]) = f(y + \delta_y \circ x).
$$
Setting $x=X$ and averaging over $X$ gives
$$
\left| f(y) - \E_X [f(y + \delta_y \circ X)] \right| =
	\left| \E_F [F(\zero)] - \E_{F,X} [F(X)] \right| \le
\E_F \left| F(\zero) - \E_{X} [F(X)] \right| \le \eps,
$$
since $F \in \FF$ with probability one and $X$ is a fractional PRG for $\FF$ with error $\eps$.
\end{proof}

\begin{claim}
\label{claim:gt_error}
Let $X_1,\ldots,X_t \in [-1,1]^n$ be independent fractional PRGs for $\FF$ with error $\eps$. Then for any $f \in \FF$,
$$
\left| f(\zero) - \E_{X_1,\ldots,X_t} [f(g_t^n(X_1,\ldots,X_t))] \right| \le t \eps.
$$
\end{claim}

\begin{proof}
The proof is by induction on $t$. The base case $t=1$ follows by definition as $g_1^n(X_1)=X_1$. For $t>1$ we will show that
$$
\left| \E[f(g_{t-1}^n(X_1,\ldots,X_{t-1}))] - \E [f(g_t^n(X_1,\ldots,X_t))] \right| \le \eps,
$$
from which the claim follows by the triangle inequality. In fact, we will show a stronger inequality: for any fixing of $x_1,\ldots,x_{t-1} \in [-1,1]^n$,
it holds that
$$
\left| f(g_{t-1}^n(x_1,\ldots,x_{t-1})) - \E_{X_t}[ f(g_t^n(x_1,\ldots,x_{t-1},X_t))] \right| \le \eps.
$$
The first inequality then follows by averaging over $x_1=X_1,\ldots,x_{t-1}=X_{t-1}$. To see why this latter inequality holds,
set $y = g_{t-1}^n(x_1,\ldots,x_{t-1})$. Then by definition,
$$
g_t^n(x_1,\ldots,x_{t-1},X_t) = y + \delta_y \circ X_t.
$$
The claim now follows from \expref{Claim}{claim:prg_around_nonzero}.
\end{proof}

We have so far proved that $g_t^n(X_1,\ldots,X_t)$ is a fractional PRG for $\FF$ with slightly worse error. Although we do not need it, it is worth
noting that it is symmetric since $X_1,\ldots,X_t$ are symmetric
and $-g_t^n(X_1,\ldots,X_t)=g_t^n(-X_1,\ldots,-X_t)$.
To conclude, we show that it
converges fast to a value close to $\{-1,1\}^n$. This is the effect of \emph{polarization}. It will be enough to analyze this
for one-dimensional random variables.

\begin{claim}
\label{claim:polar_onedim}
Let $A_1,\ldots,A_t \in [-1,1]$ be independent symmetric random variables with $\E[A_i^2] \ge p$.
For $i=1,\ldots,t$ define
$$
B_i:=g_i(A_1,\ldots,A_i) = B_{i-1} + (1-|B_{i-1}|) A_i.
$$
Then $\E[B_t^2] \ge 1-q$ where $q=3 \exp(-tp/16)$.
\end{claim}

\begin{proof}
 Let $C_i := 1-|B_i|$ be the distance to $\{-1,1\}$ at step $i$. We show that $C_i$ converges to $0$ exponentially fast. Observe that $C_i$ satisfies the following recursive definition:
$$
C_i =
\begin{cases}
C_{i-1} (1-A_i\cdot \sign(B_{i-1})) & \text{if } C_{i-1} (1-A_i\cdot\sign(B_{i-1})) \le 1, \\
2 - C_{i-1} (1-A_i\cdot \sign(B_{i-1})) & \text{if } C_{i-1} (1-A_i\cdot \sign(B_{i-1})) > 1.
\end{cases}
$$
In either case one can verify that $C_i \in [0,1]$ and that
$$
	C_i \le C_{i-1} (1-A_i\cdot\sign(B_{i-1})).
$$
	Now observe that $C_{i-1}$ and $A_i\cdot\sign(B_{i-1})$ are independent. This is because $B_{i-1}$ is symmetric (because $A_j$'s are symmetric), and so $|B_{i-1}|$ and $\sign(B_{i-1})$ are independent. So we can write, 
$$
\E\left[\sqrt{C_i}\right] \leq \E\left[\sqrt{C_{i-1}}\right] \E\left[\sqrt{1-A_i\cdot\sign(B_{i-1})}\right].
$$
The Taylor expansion of $\sqrt{1-x}$ in $[-1,1]$ is
$$
\sqrt{1-x} = 1 - \frac{x}{2} - \frac{x^2}{8}  - \frac{x^3}{16} - \ldots.
$$
	In particular, all the coefficients except for the constant term are negative. As $A_i\cdot\sign(B_{i-1})$ is symmetric, $\E[(A_i\cdot\sign(B_{i-1}))^k]=0$ for any odd $k$, so
$$
\E\left[\sqrt{1-A_i\cdot\sign(B_{i-1})}\right] \le 1 - \frac{\E[A_i^2]}{8} \le 1-\frac{p}{8} \le \exp(-p/8).
$$
Thus
$$
\E\left[\sqrt{C_t}\right] \le \prod_{i=1}^t \E\left[\sqrt{1-A_i\cdot\sign(B_{i-1})}\right] \le \exp(-t p / 8).
$$
	Now we use Markov's inequality. We know $\Pr[\sqrt{C_t}\geq \lambda \E\left[\sqrt{C_t}\right] ] \leq \lambda^{-1}$. By choosing $\lambda = \exp(t p / 16)$ we get $\Pr[C_t \geq \exp(-t p / 8) ] \leq \exp(-t p / 16)$.
 If $C_t \le \exp(-tp/8)$ then $1-B_t^2 \le 2 \exp(-tp/8)$.
If not, then we can trivially bound $1-B_t^2 \le 1$. Putting these together gives
\[
\E[1-B_t^2] \le 2 \exp(-tp/8) + \exp(-tp/16) \le 3 \exp(-tp/16).\qedhere
\]
\end{proof}

To provide a piece of intuition explaining the fast convergence of this random walk, notice that once $C_i$ becomes sufficiently small, it gets more and more difficult to increase the value of $C_i$ again. This could be best explained with an example. Suppose all $A_i$'s take value in $\{-0.5,0.5\}$. We start at $B_0 = 0$ and take a step, say $A_1 = 0.5$, and therefore $B_1 = 0.5$. Now observe that the length of the next step would be only $(1-|B_1|)|A_2| = 0.25$. So even if $A_2 = -0.5$, we get $B_2 = 0.25$, which means we still need to take one more step to become less than $0$.  In other words, once we get close to the boundary $\{-1,1\}$, the random walk  converges faster as it gets more difficult to move away from the boundary.

\begin{corollary}
\label{cor:polar}
Let $X_1,\ldots,X_t \in [-1,1]^n$ be independent symmetric $p$-noticeable random variables. Define
$Y=g_t^n(X_1,\ldots,X_t)$.
Then $Y$ is $(1-q)$-noticeable for $q=3 \exp(-tp/16)$.
\end{corollary}

\begin{proof}
Apply \expref{Claim}{claim:polar_onedim} to each coordinate of $Y$.
\end{proof}

\expref{Lemma}{lemma:amplify} follows by combining \expref{Claim}{claim:gt_error} and \Cref{cor:polar}.

\subsection{Proof of \expref{Lemma}{lemma:round}}
We prove \expref{Lemma}{lemma:round} in this section.
Let $x \in [-1,1]^n$ be a possible outcome of $X$. Let $W:=W(x) \in \{-1,1\}^n$
be a random variable, where $W_1,\ldots,W_n$ are independent and $\E[W_i]=x_i$. Then $\E_W[f(W)] = f(x)$.
As $f$ takes values in $[-1,1]$, we can upper bound $|f(x)-f(\sign(x))|$ by
$$
|f(x) - f(\sign(x))| = |\E_W [f(W)] - f(\sign(x))| \le 2\Pr[W \ne \sign(x)].
$$
The last term can be bounded by the union bound,
$$
2\Pr[W \ne \sign(x)] \le 2\sum_{i=1}^n \Pr[W_i \ne \sign(x_i)] = \sum_{i=1}^n 1-|x_i|.
$$
Setting $x=X$ and averaging over $X$ gives
$$
|\E_X [f(X)] - \E_X [f(\sign(X))]| \le
\E_X|f(X) - f(\sign(X))| \le
 \sum_{i=1}^n \E[1-|X_i|].
$$
As $X$ is $(1-q)$-noticeable it satisfies $\E[X_i^2] \ge 1-q$ for all $i$.
As $1-z \le 1-z^2$ for all $z \in [0,1]$ we have
$$
\E[1-|X_i|] \le \E[1-X_i^2] \le q.
$$
This concludes the proof as
$$
|f(\zero) - \E_X [f(\sign(X))]| \le
|f(\zero) - \E_X [f(X)]|+|\E_X [f(X)] - \E_X [f(\sign(X))]| \le \eps+qn,
$$
where the first inequality follows as $X$ is a fractional PRG with
error $\eps$, and the second by the discussion above.

\section{PRGs for functions with bounded Fourier tails}
\label{sec:spectral}

Several natural families of Boolean functions have bounded Fourier tails, such as: $\cclass{AC}^0$ circuits \cite{linial1993constant,mansour1995nlog}; functions
with bounded sensitivity \cite{gopalan2016degree,lovett2016robust}; and
functions computed by branching programs of various forms~\cite{reingold2013pseudorandomness,chattopadhyay2018improved}.
Our goal is to construct a universal PRG which fools any such function. We consider two variants: $L_1$ bounds and $L_2$ bounds.

\begin{definition}[$L_1$ bounds]
For $a,b \ge 1$, we denote by $\LL_1^n(a,b)$ the family of $n$-variate Boolean functions $f:\{-1,1\}^n \to \{-1,1\}$ which satisfy
$$
\sum_{\substack{S\subseteq [n] \\ |S| = k }} |\widehat{f}(S)| \leq a \cdot b^k \qquad \forall k=1,\ldots,n.
$$
\end{definition}

\begin{definition}[$L_2$ bounds]
For $a,b \ge 1$, we denote by $\LL_2^n(a,b)$ the family of $n$-variate Boolean functions $f:\{-1,1\}^n \to \{-1,1\}$ which satisfy
$$
\sum_{\substack{S\subseteq [n] \\ |S| \ge k }} \widehat{f}(S)^2 \leq a \cdot 2^{-k/b} \qquad \forall k=1,\ldots,n.
$$
\end{definition}

Tal~\cite{tal2017tight} showed that $L_2$ bounds imply $L_1$ bounds: if $f \in \LL_2(a,b)$ then $f \in \LL_1(a,b')$ for $b'=O(b)$.
The reverse direction is false, as can be witnessed by the PARITY function.
So, the class of functions with $L_1$ bounded Fourier tails is richer, and we focus on it.

In the following lemma, we construct a fractional PRG for this class, which we will then amplify to a PRG\@.
We note that this lemma holds also for bounded functions, not just Boolean functions.
The construction is based on a scaling of almost $d$-wise independent random variables, whose definition we now recall.

\begin{definition}[Almost $d$-wise independence]
A random variable $Z \in \{-1,1\}^n$ is $\eps$-almost $d$-wise independent if,
for any restriction of $Z$ to $d$ coordinates,
the marginal distribution has statistical distance at
most $\eps$ from the uniform distribution on $\{-1,1\}^d$.
\end{definition}

Naor and Naor \cite{naor1993small} gave an explicit construction of an $\eps$-almost $d$-wise random variable
$Z \in \{-1,1\}^n$ with seed length $O(\log\log n +  d + \log(1/\eps))$. We note that this seed length is optimal, up to the hidden constants.

\begin{lemma}\label{lemma:L1}
Fix $n,a,b \ge 1$ and $\eps>0$. There exists a fractional PRG $X \in [-1,1]^n$ that fools $\LL_1^n(a,b)$
with error $\eps$, such that
\begin{enumerate}[(i)]
\item $X$ is $p$-noticeable for $p={1}/{(4 b^2)}$.
\item The seed length of $X$ is $O(\log \log n + \log(a/\eps))$.
\end{enumerate}
\end{lemma}

\begin{proof}
Fix $f \in \LL_1^n(a,b)$.
Set $d=\lceil \log 2a/ \eps \rceil, \delta=\eps/2a, \beta=1/2b$. Let $Z \in \{-1,1\}^n$ be a $\delta$-almost $d$-wise independent random variable,
and set $X=\beta Z$ which takes values in $\{-\beta,\beta\}^n$. We claim that $X$ satisfies the requirements of the lemma. Claim (i) clearly holds, and claim (ii) holds by the Naor-Naor construction. We thus focus on proving that $X$ fools $\FF$ with error $\eps$.

Fix $f \in \FF$ and consider its Fourier expansion:
$$
f(x) = \sum_{S \subseteq [n]} \widehat{f}(S)  x^S.
$$
We need to show that $\E[f(X)]$ is close to $f(\zero)$. Averaging over $X$ gives
$$
|\E[f(X)]-f(\zero)| \le \sum_{|S|>0} |\widehat{f}(S)| \cdot |\E[X^S]| =
\sum_{|S|>0} |\widehat{f}(S)| \cdot \beta^{|S|} |\E[Z^S]|.
$$
We next bound $|\E[Z^S]|$. If $|S| \le d$ then by the definition of $Z$
we have $|\E[Z^S]| \le \delta$. If $|S|>d$ we bound trivially $|\E[Z^S]| \le 1$. Let $W_k = \sum_{S: |S|=k } |\hat{f}(S)|$, where by assumption $W_k \le a \cdot b^k$. Thus
$$
|\E[f(X)]-f(\zero)| \le \delta \sum_{k=1}^d W_k \beta^k + \sum_{k>d} W_k \beta^k \le  \delta a \sum_{k=1}^d (\beta b)^k + a \sum_{k>d} (\beta b)^k \le \delta a + 2^{-d} a
$$
where we used the choice of $\beta = 1/2b$. The claim follows as we set $\delta=\eps/2a $ and $2^{-d} \le \eps/2a$.
\end{proof}

Applying \expref{Theorem}{thm:main} using the fractional PRG constructed in \expref{Lemma}{lemma:L1} gives the following PRG construction. Note that
we still need to require that $\FF$ is closed under restrictions.

\begin{theorem}
\label{thm:L1}
Let $\FF$ be a family of $n$-variate Boolean functions closed under restrictions. Assume that $\FF \subset \LL_1^n(a,b)$ or that $\FF \subset \LL_2^n(a,b)$.
	Then, for any $\eps\leq {1}/{\poly(b \log n)}$ there exists an explicit PRG $X \in \{-1,1\}^n$ which fools $\FF$ with error $\eps>0$,
whose seed length is $O(\log(n/\eps)  (\log\log n + \log(a/\eps)) b^2)$.
\end{theorem}

\subsection{Applications}
We apply our PRG from \expref{Theorem}{thm:L1} to several well studied classes of Boolean functions that are known to satisfy a Fourier tail bound.

\subsubsection{Functions of bounded sensitivity}
Let $f:\{-1,1\}^n \to \{-1,1\}$ be a Boolean function. Its sensitivity at an input $x \in \{-1,1\}^n$ is the number of neighbors $x'$ of $x$
(that is, $x'$ and $x$ differ at exactly one coordinate)
such that $f(x') \ne f(x)$. The (max) sensitivity of $f$ is $s(f)=\max_x s(f,x)$. The sensitivity conjecture speculates
that functions of sensitivity $s$ can be computed by decision trees of depth $\poly(s)$. A corollary would be that almost $\poly(s)$-wise
distributions fool functions of low sensitivity. So, one may ask to construct comparable PRGs for functions of low sensitivity.

This question was first considered by Hatami and Tal~\cite{hatami2017pseudorandom}. They constructed a PRG with subexponential seed length
$\exp(O(\sqrt{s}))$. \expref{Theorem}{thm:L1} gives an improved construction that essentially matches the consequence of the sensitivity conjecture. Our PRG uses the recent bounds of Gopalan {et al.}~\cite{gopalan2016degree}
on the Fourier tail of functions of low sensitivity. Concretely,
Gopalan {et al.}~\cite{gopalan2016degree} show that if $s(f)=s$ then $f \in \LL_1(1,t)$ for $t=O(s)$.
It is straightforward to verify that a restriction can only decrease the sensitivity of the function, so that the
class of functions of sensitivity at most $s$ is closed under restrictions.
A direct application of \expref{Theorem}{thm:L1} gives a PRG with seed length $O(s^2 \log(n/\eps) (\log\log(n) + \log(1/\eps)))$.

To get a somewhat improved bound, one can apply a result of Simon~\cite{simon1983tight} that shows that if $s(f)=s$
then $f$ depends on at most $m=2^{O(s)}$ many inputs. In this case, the analysis of \expref{Theorem}{thm:main} can be applied
with $m$ variables instead of $n$ variables , so that we only need $O(\log m/\eps)$ iterations. Note that the fractional PRG still requires
a seed length which depends on the original $n$. We obtain:

\begin{corollary}
\label{cor:sensitivity}
For any $n,s \ge 1$ and $\eps\leq 1/\poly(s)$, there exists an
explicit PRG which fools $n$-variate Boolean functions with
sensitivity $s$ with error $\eps$, whose seed length is
\[
  O\bigl(s^3 \log\bigl(1/\eps) (\log\log n + \log(1/\eps))\bigr).
\]
\end{corollary}

We note that the $\log\log n$ term cannot be removed. Indeed, even if we restrict attention to functions which are XOR of at most $2$ bits (for which $s=2$) the seed length required is $\Omega(\log\log n + \log(1/\eps))$.

\subsubsection{Unordered branching programs}
An oblivious read-once branching program (abbrv ROBP) $B$ of width $w$ is a non-uniform model of computation, that captures randomized algorithms with space $\log w$.  A branching program $B$ maintains a state in the set $\{1,\ldots,w\}$ and reads the input bits in a known fixed order. At time step $i =1,\ldots,n$, $B$ reads a bit and based on the time step, the read bit and the current state it transitions to a new state.
Thus, $B$ can be thought of as a layered directed graph, with $w$ nodes in each layer, and two edges going out of each node to the immediately next layer, one labeled with a $1$ and the other labeled with a $-1$.

Let $\BB^n(w)$ be the class of $n$-variate Boolean functions computed by read-once oblivious branching programs of width $w$,
where the order of the inputs is arbitrary.
A recent work of Chattopadhyay {et al.}~\cite{chattopadhyay2018improved} showed that these functions
have $L_1$ bounded Fourier tails. Concretely, $\BB^n(w) \subset \LL_1^n(1,t)$ for $t=O((\log n)^w)$.
They used this to construct a PRG with seed length $O(\log n)^{w-1} \log^2(n/\epsilon)\log \log n$.
Using our PRG from \expref{Theorem}{thm:L1} we get a comparable (although slightly worse) seed length.
Note that $\BB^n(w)$ is closed under restrictions.

\begin{corollary}
\label{cor:ROBP}
	Fix $n,w \ge 1$ and $\eps\leq 1/(\log n)^{O(w)}$. There is an explicit PRG which fools $\BB^n(w)$ with error $\eps>0$,
whose seed length is $O(\log (n/\eps) (\log \log n + \log 1/\eps) (\log n)^{2w})$.
\end{corollary}


\subsubsection{Permutation branching programs}
A special case of read-once branching programs are permutation branching programs, where the transition function from level $i$
to level $i+1$ in the graph is a permutation for every choice of the input bit. We denote it by $\BBperm^n(w) \subset \BB^n(w)$.
Reingold {et al.}~\cite{reingold2013pseudorandomness} showed that if a Boolean function
is computed by a permutation branching program of width $w$, then it has $L_2$ bounded Fourier tails with parameter $2 w^2$.
Note that permutation branching programs are also closed under restrictions. Thus we obtain the following result:

\begin{corollary}
\label{cor:perm_BP}
	Fix $n,w \ge 1$ and $\eps\leq 1/(\poly(w, \log n))$. There is an explicit PRG which fools $\BBperm^n(w)$ with error $\eps>0$,
whose seed length is $O(\log (n/\eps) (\log \log n + \log 1/\eps) w^4)$.
\end{corollary}
%
%
The dependence on $n$ in our PRG is better than the result of Reingold et al.~\cite{reingold2013pseudorandomness},
as they obtained seed length $O(w^2 \log(w) \log(n) \log(nw/\eps) + w^4 \log^2(w/\eps))$.

%
%
Reingold et al.~\cite{reingold2013pseudorandomness} actually show the Fourier tail bounds for a more general
class of branching programs, called regular branching programs. However, these are not closed under restriction,
%
%
and hence our PRG construction fails to work (the same problem occurs also in the construction from~\cite{reingold2013pseudorandomness}).

%
\subsubsection{Bounded-depth circuits}   %
The class of bounded-depth Boolean circuits $\cclass{AC}^0$ has been widely studied. In particular, Linial, Mansour and Nisan~\cite{linial1993constant} showed that it has bounded $L_2$ Fourier tails. Tal~\cite{tal2017tight} obtained improved
bounds. If $f$ is an $n$-variate Boolean function computed by an $\cclass{AC}^0$ circuit of depth $d$ and size $s$,
then $f \in \LL_2(n,t)$ for $t=2^{O(d)}\log^{d-1}s$. \expref{Theorem}{thm:L1} provides a new PRG for $\cclass{AC}^0$ which is comparable
with the existing PRGs of Nisan~\cite{nisan1991pseudorandom}, Braverman~\cite{braverman2010polylogarithmic}, Tal~\cite{tal2017tight}, and Trevisan and Xue~\cite{trevisan2013derandomized}.

\begin{corollary}
\label{cor:AC0}
	Fix $n,s \ge 1$ and $\eps\leq 1/(\log s)^{O(d)}(\log n)^{O(1)}$. There is an explicit PRG which fools $n$-variate functions which can be computed
by $\cclass{AC}^0$ circuits of size $s$ and depth $d$, with error $\eps>0$,
whose seed length is $O(\log (n/\eps) (\log \log n + \log 1/\eps) \log^{2d-2} s)$.
\end{corollary}

\section{PRG for functions which simplify under random restriction}
\label{sec:random_restriction}
Another generic application of our framework is constructing PRGs for classes that simplify under random restriction.
Let $\FF$ be a family of functions  $f:\{-1,1\}^n \to \{-1,1\}$  which are extended multilinearly to $[-1,1]^n$. Fix a parameter $0<p<1$ and define the $p$-averaged function of $f$, denoted $f_p:\{-1,1\}^n \to [-1,1]$ as follows:
sample $A \subseteq [n]$ where $\Pr[i \in A]=p$ independently for $i \in [n]$, and define
$$f_p(x) = \E_A\E_U [f(x_A,U_{A^c})]$$
where $x_A \in \{-1,1\}^A$ is the restriction of the input $x$ to the coordinates in $A$, and $U \in \{-1,1\}^n$ is independently and uniformly chosen.
 \begin{claim}
  \label{claim:p_averaged}
 $f_p (x) = f(px)$.
 \end{claim}
 \begin{proof}
 Let $A,U$ be random variables as defined above. Define a random variable
 $Y \in \{-1,1\}^n$ as follows:
$$
Y_i =
\begin{cases}
x_i & \text{if } A_i=1,\\
U_i & \text{if } A_i=0.
\end{cases}
$$
Note that $Y$ is a product distribution.
By definition of $f_p$, $f_p(x) = \E[f(Y)]$. By multilinearity of $f$,
$\E[f(Y)]=f(\E[Y]) = f(px)$.
\end{proof}
Overall, using \expref{Claim}{claim:p_averaged}, we obtain the following corollary of \expref{Theorem}{thm:main}.
\begin{corollary}
%
%
	
Suppose that we have a standard PRG $X$ for the class of $p$-averaged functions $\FF_p = \{f_p : f\in \FF\}$ with error $\epsilon$. Then $X'=pX$ is a fractional PRG for the class $\FF$. Therefore, we get a standard PRG for $\FF$ using $t = O(\log(n/\epsilon)/p^2)$ independent copies of $X$ that has error $t\epsilon$.
\end{corollary}

\section{Spectral tail bounds for low-degree \texorpdfstring{$\F_2$}{F2}-polynomials}
\label{section:lowdegree_spectral}
In this section, we prove  $L_1$ Fourier tail bounds for functions computed by low-degree polynomials on $\F_2$. However, our bounds fall short of implying PRGs for the class of low-degree $\F_2$ polynomials in our framework.

\begin{theorem}
\label{thm:polynomialspectralnorm} Let $p:\F_2^n \to \F_2$ be a polynomial
of degree $d$, and let $f(x)=(-1)^{p(x)}$. Then
$$
\sum_{\substack{S\subseteq [n] \\ |S| = k }} |\widehat{f}(S)| \leq   (k2^{3d})^k \qquad \forall k=1,\ldots,n.
$$
\end{theorem}

We note that $L_2$ bounds do not hold for low-degree polynomials, as can be witnessed by taking a high-rank quadratic polynomial. We prove \expref{Theorem}{thm:polynomialspectralnorm} in the remainder of this section.
For convenience, with a slight abuse of notation, we use the basis $f:\{-1,1\}\rightarrow \{-1,1\}$.
We first introduce some notation to simplify the presentation. Let
$$
W_k(f) := \sum_{|S|=k} |\hat{f}(S)|
$$
denote the weight of the level-$k$ Fourier coefficients of a Boolean
function $f$, and let
$$
W(d,k) := \max \{W_k(f): f=(-1)^p, \; \deg(p) \le d\}
$$
be the maximum of $W_k$ over degree-$d$ polynomials. Note that we do not make any assumption on the number of variables $n$. We prove the following lemma from which \expref{Theorem}{thm:polynomialspectralnorm} follows relatively easily.
\begin{lemma}
\label{lemma:rec_poly} For any $d,k \ge 1$, $$ W(d,k)^{2} \le 2^{2k}W(d-1,2k) + W(d,k) \cdot \sum_{\ell=1}^{k} \binom{k}{\ell}   W(d, k -\ell).$$
\end{lemma}

We first show that \expref{Theorem}{thm:polynomialspectralnorm} follows easily from \expref{Lemma}{lemma:rec_poly}.

\begin{proof}[Proof of \expref{Theorem}{thm:polynomialspectralnorm} given \expref{Lemma}{lemma:rec_poly}]
The proof of \expref{Theorem}{thm:polynomialspectralnorm} is by induction, first on $d$ and then on $k$. The base case of $d=1$ is straightforward, so assume $d \ge 2$. By \expref{Lemma}{lemma:rec_poly} we have
\begin{align*}
W(d,k)^2 &\le 2^{2k} \left( 2k \cdot 2^{3(d-1)} \right)^{2k}+W(d,k)
\sum_{\ell=1}^k \binom{k}{\ell} \left((k-\ell) 2^{3d} \right)^{k-\ell} \\
& \le
\left( k \cdot 2^{3d-1} \right)^{2k} + W(d,k) \sum_{\ell=1}^k \binom{k}{\ell} \left((k-1) 2^{3d} \right)^{k-\ell} \\
& = \left( k \cdot 2^{3d-1} \right)^{2k} + W(d,k) \left(\left((k-1) 2^{3d}+1\right)^k - \left((k-1) 2^{3d}\right)^k \right).
\end{align*}
Assume towards a contradiction that $W(d,k)>(k 2^{3d})^k$. Dividing by $W(d,k)$ on both sides gives
\begin{align*}
W(d,k) & \le \left( k \cdot 2^{3d-1} \right)^{k} + \left((k-1) 2^{3d}+1\right)^k - \left((k-1) 2^{3d}\right)^k.
\end{align*}
If $k=1$ then we reach a contradiction as $2^{3d-1}+1 \le 2^{3d}$. If $k>1$ then as $(k-1) 2^{3d} \ge k 2^{3d-1}$
the first term gets canceled by the third term, and the second term is
at most $(k 2^{3d})^k$. In either case, we reached a contradiction.
\end{proof}

{}From now on we focus on proving \expref{Lemma}{lemma:rec_poly}.
To that end, fix $f$ computed by a polynomial
of degree $d$ which maximizes $W_k(f)$. We shorthand $g(S) = |\hat{f}(S)|$.
The following claims are used in the proof of \expref{Lemma}{lemma:rec_poly}.
\begin{claim}
\label{claim:fourier_level_down}
For any $0\le a < b \le n$ and $A \subset [n]$ of size $|A| = a$,
$$\sum_{B: |B|=b, A \subseteq B} g(B) \le W(d,b -a).$$
\end{claim}
\begin{claim}
\label{claim:derivative}
$\displaystyle \sum_{S,T: |S|=|T|=k, S \cap T = \emptyset} g(S)g(T) \le 2^{2k} W(d-1,2k)$.
\end{claim}
\begin{claim}
\label{claim:sublevels} For any $1 \le \ell \le k$, $$ \sum_{S,T: |S|=|T|=k, |S \cap T| = \ell} g(S) g(T) \le \binom{k}{\ell} W(d,k)W(d,k-\ell). $$
\end{claim}

We first show how to prove \expref{Lemma}{lemma:rec_poly} using the above claims.
\begin{proof}[Proof of \expref{Lemma}{lemma:rec_poly}]
We have,
\begin{align*}
W(d,k)^2 &= \sum_{S,T: |S|=|T|=k} g(S) g(T) \\
 &= \sum_{\ell =0}^{k} \sum_{S, T: |S|= |T|=k, |S \cap T|=\ell} g(S) g(T) \\
 &= \sum_{S,T: |S|= |T|=k, S \cap T = \emptyset} g(S) g(T) + \sum_{\ell =1}^{k} \sum_{S, T: |S|= |T|=k, |S \cap T|=\ell} g(S) g(T) \\
 &\le 2^{2k}W(d-1,2k) + W(d,k) \cdot \sum_{\ell=1}^{k} \binom{k}{\ell}  W(d, k -\ell),
\end{align*} where the last inequality follows by using the bounds from \expref{Claim}{claim:derivative} and \expref{Claim}{claim:sublevels}.
\end{proof}
We now proceed to prove the missing claims.
\begin{proof}[Proof of \expref{Claim}{claim:fourier_level_down}] We use induction on $a$ and $b$. The claim is direct for $a =0$  and any $b > a$. Thus suppose $b>a>0$ and let $ i\in A$. Let $A' = A \setminus \{ i\}$. We have
\begin{align*}
\sum_{B:|B|=k, A \subseteq B} g(B) &=  \sum_{B' \subseteq [n] \setminus \{ i\}:|B'|=b-1, A' \subseteq B' } g(B' \cup \{ i\}) \\
& = \sum_{B' \subseteq [n] \setminus \{ i\}:|B'|=b-1, A' \subseteq B' } |\widehat{f}(B' \cup \{ i\})|.
\end{align*}
Let $f_{i\rightarrow1}$ and $f_{i\rightarrow-1}$ be the functions obtained from $f$ by setting the $i$'th bit to $1$ and $-1$, respectively. It is easy to verify that $|\widehat{f}(B \cup \{ i\})| = \frac{1}{2}|\widehat{f_{i\rightarrow1}}(B) - \widehat{f_{i\rightarrow-1}}(B)|$. Thus, continuing with our estimate, we have
\begin{align*}
\sum_{B:|B|=b, A \subseteq B} g(B)  &\le \frac{1}{2}\sum_{B':|B'|=b-1, A' \subseteq B' \cup \{ i\} } \big(|\widehat{f_{i\rightarrow1}}(B')| + |\widehat{f_{i\rightarrow-1}}(B')|\big) \\
&\le W(d,(b-1)- (a-1))=W(d,b-a),
\end{align*} where the last inequality follows from induction hypothesis.
\end{proof}
\begin{proof}[Proof of \expref{Claim}{claim:derivative}] For any $S \subseteq [n]$, let $e_S \in \{-1,1\}$ be the sign of $\widehat{f}(S)$, so that $g(S) = e_S \cdot \widehat{f}(S)$. Let $X,Y,Z$ be independent uniform distributions on $\{ -1,1\}^n$. We have
\begin{align*}
\sum_{S,T: |S|=|T|=k, S \cap T = \emptyset} g(S) g(T)
&= \sum_{S,T: |S|=|T|=k, S \cap T = \emptyset} e_S e_T \E_Y[f(Y)Y^S] \cdot \E_Z[f(Z)Z^T] \\
&= \sum_{S,T: |S|=|T|=k, S \cap T = \emptyset} e_S e_T \E_{Y,Z}[f(Y) Y^S f(Z) Z^T] \\
&= \sum_{S,T: |S|=|T|=k, S \cap T = \emptyset} e_S e_T \E_{X,Y,Z}[f(X\circ Y)f(X \circ Z) X^{S \cup T} Y^S  Z^T].
\end{align*}
This follows as $(Y,Z)$ and $(X \circ Y,X \circ Z)$ are identically distributed.
	Now consider any fixing of $Y=y$ and $Z=z$. Define the function $h_{y,z}(x) = f(x\circ y)f(x \circ z)$. Recall that $f$ is a degree-$d$ polynomial. Thus $h_{y,z}(x)= f(x) f(x\circ z\circ y)$ is the multiplicative derivative
	of $f$ in direction $y \circ z$. In particular, the degree of $h_{y,z}$ is at most $d-1$, since taking
%
derivatives   %
always reduces
the    %
degree.  More information about properties of derivatives can be found
in \cite{o2014analysis}.
Thus we have
	\begin{align*}\sum_{S,T: |S|=|T|=k, S \cap T = \emptyset} e_S e_T y^S z^T \E[f(X \circ y)f(X \circ z)X^{S \cup T}] &\le \binom{2k}{k} \sum_{R: |R|=2k}\Big|\E[h_{y,z}(X)X^{R}]\Big| \\
& \le 2^{2k} W(d-1,2k).
\end{align*}
The proof follows now by noting that the above bound holds for any choice of $y$ and $z$, and then averaging over $y=Y,z=Z$.
\end{proof}

\begin{proof}[Proof of \expref{Claim}{claim:sublevels}] We have,
\begin{align*}
\sum_{S,T: |S|=|T|=k, |S \cap T| = \ell} g(S) g(T) & \le \sum_{L: |L| = \ell} \Big( \sum_{S: |S|= k, L \subseteq S} g(S)\Big)^2 \\
& \le \Big( \max_{L: |L|= \ell} \sum_{S: |S|= k, L \subseteq S} g(S) \Big) \Big(\sum_{L,S:|L|=\ell, |S|=k, L \subseteq S} g(S) \Big) \\
&\le W(d,k-\ell) \cdot \Big(\sum_{S: |S|=k} \sum_{L: L\subseteq S, |L|= \ell} g(S) \Big) &&\text{(using \expref{Claim}{claim:fourier_level_down})} \\
&\le W(d,k-\ell) \cdot  \binom{k}{\ell}\cdot  W(d,k). &&\qedhere
\end{align*}
\end{proof}


\section{Smoothness}
\label{section:smoothness}
In this section we provide a partial answer for
\expref{Open Problem}{open:early_termination}, regarding early termination of the random walk. Let $Y_t \in [-1,1]^n$ be the location of the random walk at time $t$. We would like to guarantee that if  $Y_t$ is close enough to $\sign(Y_t)$ then we can round $Y_t$ to $\sign(Y_t)$ without changing the value of $f$ by too much. Therefore, given $f:[-1,1]^n\to [-1,1]$, it would be desirable to show $f$ is ``smooth'' enough: there is a bound $W$ such that
$$
\forall \alpha,\beta \in [-1,1]^n, \ |f(\alpha)-f(\beta)| \le W \|\alpha-\beta\|_\infty.
$$
Observe that should such a $W$ exist, then if at some step $t$ we have $\|Y_t-\sign(Y_t)\|_\infty\leq \eps/W$ ,  then we can terminate the random walk immediately and guarantee that $\|f(Y_t)-f(\sign(Y_t))\|_\infty\leq \eps$.
We show that such smoothness property holds for functions with bounded sensitivity.


\subsection{Bounded sensitivity functions}

We show that smoothness follows from a bound on the (maximum) sensitivity of a Boolean function.

\begin{lemma}
\label{lemma:sensitivity_smoothness}
Let $f:\{-1,1\}^n\to [-1,1]$ be a boolean function with maximum sensitivity $s$. Then, for any $\alpha,\beta\in [-1,1]^n$ it holds that
$$
|f(\alpha) - f(\beta) | \leq 8s \|\alpha -\beta\|_\infty.
$$
\end{lemma}

We first consider the case that $\|\alpha-\beta\|_\infty$ is very small.
\begin{claim}
\label{claim:small_distance}
Let $f:\{-1,1\}^n\to [-1,1]$ be a Boolean function with maximum sensitivity $s$. 
Let $\alpha,\beta\in [-1,1]^n$ such that $\|\alpha-\beta\|_\infty \leq 1/n^2$. Then
$$
|f(\alpha) - f(\beta) | \leq 8s \|\alpha -\beta\|_\infty.
$$
\end{claim}


To prove the result for arbitrary $\alpha,\beta\in [-1,1]^n$ using \expref{Claim}{claim:small_distance}, consider the line segment from $\alpha$ to $\beta$ and integrate $f$ along that line segment. Thus, \expref{Lemma}{lemma:sensitivity_smoothness} follows directly from \expref{Claim}{claim:small_distance}.

\begin{proof}[Proof of \expref{Claim}{claim:small_distance}]
Let $\delta=\|\alpha-\beta\|_{\infty}$.
We first consider the easier case of $\alpha\in \{-1,1\}^n$.
Pick $b\in\{-1,1\}^n$ randomly by flipping each coordinate of $\alpha$ independently with probability $|\alpha_i-\beta_i|/2$ so that $\E f(b) = f(\beta)$. Note that $f(b) \neq f(\alpha)$ if either exactly one sensitive coordinate of $\alpha$ is flipped, which occurs with probability at most $s \delta$, or if at least two coordinates get flipped, which occurs with probability at most $(n \delta)^2$. Therefore
$$
|f(\alpha) - \E f(b) | \leq s \delta + n^2 \delta^2 \leq 2s \delta
$$
given our assumption on $\delta$.

Next, consider the general case of $\alpha \in [-1,1]^n$.
This case requires introducing an extra point $\gamma\in [-1,1]^n$ in a way that allows us to prove
\begin{equation}
\label{equation:alpha}
|f(\alpha) - f(\gamma)| \leq 4s\cdot \|\alpha -\gamma\|_\infty
\end{equation}
and
\begin{equation}
\label{equation:beta}
|f(\beta) - f(\gamma)| \leq 4s\cdot \|\beta -\gamma\|_\infty
\end{equation}
separately.
We choose $\gamma$ in a way that $\forall i \in [n], \gamma_i = \alpha_i$ or $\gamma_i = \beta_i$. These equations altogether give the claim.
To choose $\gamma$, let $S\subseteq [n]$ be the set of coordinates that $|\alpha_i| < |\beta_i|$ and pick $\gamma_i = \alpha_i$ if $i\in S$, and $\gamma_i = \beta_i$ otherwise.

We next prove \expeqref{Equation}{equation:alpha}.  The proof of \expeqref{Equation}{equation:beta} is analogous. Consider a joint random variable $(a,c)$ that satisfies the following properties:
\begin{enumerate}
\item  $a\in\{-1,1\}^n, c\in[-1,1]^n$, $\E a = \alpha$, and $\E c = \gamma$.
\item The marginal distributions of $a$ and $c$ are product distributions.
\item $\|a - \E_c[  c|a]\|_\infty \leq 2 \|\alpha-\gamma\|_\infty$ holds with probability one.
\end{enumerate}

Observe that given such $(a,c)$,
$$
|f(\alpha) - f(\gamma)| = |\E_{a,c} [f(a) - f(c)]| \le \E_{a} \left|f(a) - \E_c[f(c)|a] \right|\leq 4s\cdot \|\alpha -\gamma\|_\infty,
$$
where the last inequality uses the first case in the proof, as $a \in \{-1,1\}^n$.

Now let us construct the joint random variable $(a,c)$.
Fix  $i\in[n]$ and suppose without loss of generality that $\alpha_i \ge 0$.
Note that by construction $-\alpha_i \le \gamma_i \leq \alpha_i$. First sample $a_i$ so that $\E[a_i]=\alpha_i$.
If $a_i = -1$ then set $c_i=-1$, otherwise set
\[
  c_i = \frac{2\gamma_i + 1-\alpha_i}{1+\alpha_i}.
\]
It's easy to check that this choice of $(a,c)$ satisfies the required
conditions, finishing the proof.
\end{proof}















\section*{Acknowledgments}
We thank Avi Wigderson and David Zuckerman for various stimulating discussions during the course of our work.

%% Theory of Computing
%% File: v015a010.bbl
%% Publication date: October 21, 2019
%% Authors: Eshan Chattopadhyay, Pooya Hatami, Kaave Hosseini, and Shachar Lovett
% Define empty bibhead if not already defined
\providecommand{\bibhead}[1]{}
% Define tocrefpdfbookmark if not already defined
\expandafter\ifx\csname pdfbookmark\endcsname\relax%
  \providecommand{\tocrefpdfbookmark}{}
\else\providecommand{\tocrefpdfbookmark}{%
   \hypertarget{tocreferences}{}%
   \pdfbookmark[1]{References}{tocreferences}}%
\fi

\tocrefpdfbookmark
\begin{thebibliography}{10}

\bibitem{ajtai1985deterministic}\bibhead{ajtai1985deterministic}
{\sc Mikl{\'o}s Ajtai and Avi Wigderson}: Deterministic simulation of
  probabilistic constant depth circuits.
\newblock {\em Advances in Computing Research}, 5:199--222, 1989.
\newblock Preliminary version in
  \href{https://doi.org/10.1109/SFCS.1985.19}{FOCS'85}.

\bibitem{alon1992simple}\bibhead{alon1992simple}
{\sc Noga Alon, Oded Goldreich, Johan H{\aa}stad, and Ren{\'e} Peralta}: Simple
  constructions of almost $k$-wise independent random variables.
\newblock {\em Random Structures \& Algorithms}, 3(3):289--304, 1992.
\newblock Preliminary version in
  \href{https://ieeexplore.ieee.org/document/89575}{FOCS'90}.
\newblock [\epfmtdoi{10.1002/rsa.3240030308}]

\bibitem{bansal2010constructive}\bibhead{bansal2010constructive}
{\sc Nikhil Bansal}: Constructive algorithms for discrepancy minimization.
\newblock In {\em Proc. 51st FOCS}, pp. 3--10. IEEE Comp. Soc. Press, 2010.
\newblock [\epfmtdoi{10.1109/FOCS.2010.7}, \epfmt{arxiv}{1002.2259}]

\bibitem{braverman2010polylogarithmic}\bibhead{braverman2010polylogarithmic}
{\sc Mark Braverman}: Polylogarithmic independence fools ${AC}^0$ circuits.
\newblock {\em J. ACM}, 57(5):28, 2010.
\newblock Preliminary version in
  \href{https://ieeexplore.ieee.org/document/5231177}{CCC'09}.
\newblock [\epfmtdoi{10.1145/1754399.1754401}]

\bibitem{conf-version}\bibhead{conf-version}
{\sc Eshan Chattopadhyay, Pooya Hatami, Kaave Hosseini, and Shachar Lovett}:
  Pseudorandom generators from polarizing random walks.
\newblock In {\em Proc. 33rd Computational Complexity Conf. (CCC'18)}, pp.
  1:1--1:21. Leibniz-Zentrum fuer Informatik, 2017.
\newblock [\epfmtdoi{10.4230/LIPIcs.CCC.2018.1}]

\bibitem{chattopadhyay2018improved}\bibhead{chattopadhyay2018improved}
{\sc Eshan Chattopadhyay, Pooya Hatami, Omer Reingold, and Avishay Tal}:
  Improved pseudorandomness for unordered branching programs through local
  monotonicity.
\newblock In {\em Proc. 50th STOC}, pp. 363--375. ACM Press, 2018.
\newblock [\epfmtdoi{10.1145/3188745.3188800}]

\bibitem{goldreich2010primer}\bibhead{goldreich2010primer}
{\sc Oded Goldreich}: {\em A Primer on Pseudorandom Generators}.
\newblock Volume~55 of {\em Univ. Lecture Ser.}
\newblock Amer. Math. Soc., 2010.
\newblock
  \href{http://www.wisdom.weizmann.ac.il/~oded/prg-primer.html}{author's home
  page}.

\bibitem{gopalan2012better}\bibhead{gopalan2012better}
{\sc Parikshit Gopalan, Raghu Meka, Omer Reingold, Luca Trevisan, and Salil
  Vadhan}: Better pseudorandom generators from milder pseudorandom
  restrictions.
\newblock In {\em Proc. 53rd FOCS}, pp. 120--129. IEEE Comp. Soc. Press, 2012.
\newblock [\epfmtdoi{10.1109/FOCS.2012.77}, \epfmt{arxiv}{1210.0049}]

\bibitem{gopalan2016degree}\bibhead{gopalan2016degree}
{\sc Parikshit Gopalan, Rocco~A Servedio, and Avi Wigderson}: Degree and
  sensitivity: tails of two distributions.
\newblock In {\em Proc. 31st Computational Complexity Conf. (CCC'16)}, pp.
  13:1--13:23. Schloss Dagstuhl--Leibniz-Zentrum fuer Informatik, 2016.
\newblock [\epfmtdoi{10.4230/LIPIcs.CCC.2016.13}, \epfmt{arxiv}{1604.07432}]

\bibitem{hatami2017pseudorandom}\bibhead{hatami2017pseudorandom}
{\sc Pooya Hatami and Avishay Tal}: Pseudorandom generators for low sensitivity
  functions.
\newblock In {\em Proc. 9th Innovations in Theoretical Computer Sci. Conf.
  (ITCS'18)}, pp. 29:1--29:13. Leibniz-Zentrum fuer Informatik, 2018.
\newblock [\epfmtdoi{10.4230/LIPIcs.ITCS.2018.29}]

\bibitem{hoory2006expander}\bibhead{hoory2006expander}
{\sc Shlomo Hoory, Nathan Linial, and Avi Wigderson}: Expander graphs and their
  applications.
\newblock {\em Bull. Amer. Math. Soc.}, 43(4):439--561, 2006.
\newblock [\epfmtdoi{10.1090/S0273-0979-06-01126-8}]

\bibitem{linial1993constant}\bibhead{linial1993constant}
{\sc Nathan Linial, Yishay Mansour, and Noam Nisan}: Constant depth circuits,
  {F}ourier transform, and learnability.
\newblock {\em J. ACM}, 40(3):607--620, 1993.
\newblock Preliminary version in
  \href{https://ieeexplore.ieee.org/document/63537}{FOCS'89}.
\newblock [\epfmtdoi{10.1145/174130.174138}]

\bibitem{lovett2012constructive}\bibhead{lovett2012constructive}
{\sc Shachar Lovett and Raghu Meka}: Constructive discrepancy minimization by
  walking on the edges.
\newblock {\em SIAM J. Comput.}, 44(5):1573--1582, 2015.
\newblock Preliminary version in
  \href{https://ieeexplore.ieee.org/document/6375282}{FOCS'12}.
\newblock [\epfmtdoi{10.1137/130929400}, \epfmt{arxiv}{1203.5747}]

\bibitem{lovett2016robust}\bibhead{lovett2016robust}
{\sc Shachar Lovett, Avishay Tal, and Jiapeng Zhang}: The robust sensitivity of
  {B}oolean functions.
\newblock In {\em Proc. 29th ACM-SIAM Symp. on Discrete Algorithms (SODA'18)},
  pp. 1822--1833. SIAM, 2018.
\newblock [\epfmtdoi{10.1137/1.9781611975031.119}]

\bibitem{mansour1995nlog}\bibhead{mansour1995nlog}
{\sc Yishay Mansour}: An $n^{O(\log \log n)}$ learning algorithm for {DNF}
  under the uniform distribution.
\newblock {\em J. Comput. System Sci.}, 50(3):543--550, 1995.
\newblock Preliminary version in
  \href{https://dl.acm.org/citation.cfm?doid=130385.130391}{COLT'92}.
\newblock [\epfmtdoi{10.1006/jcss.1995.1043}]

\bibitem{naor1993small}\bibhead{naor1993small}
{\sc Joseph Naor and Moni Naor}: Small-bias probability spaces: {E}fficient
  constructions and applications.
\newblock {\em SIAM J. Comput.}, 22(4):838--856, 1993.
\newblock Preliminary version in
  \href{https://dl.acm.org/citation.cfm?id=100244}{STOC'90}.
\newblock [\epfmtdoi{10.1137/0222053}]

\bibitem{nisan1991pseudorandom}\bibhead{nisan1991pseudorandom}
{\sc Noam Nisan}: Pseudorandom bits for constant depth circuits.
\newblock {\em Combinatorica}, 11(1):63--70, 1991.
\newblock [\epfmtdoi{10.1007/BF01375474}]

\bibitem{nisan1988hardness}\bibhead{nisan1988hardness}
{\sc Noam Nisan and Avi Wigderson}: Hardness vs randomness.
\newblock {\em J. Comput. System Sci.}, 49(2):149--167, 1994.
\newblock Preliminary version in
  \href{https://ieeexplore.ieee.org/document/21916}{FOCS'88}.
\newblock [\epfmtdoi{10.1016/S0022-0000(05)80043-1}]

\bibitem{o2014analysis}\bibhead{o2014analysis}
{\sc Ryan O'Donnell}: {\em Analysis of Boolean Functions}.
\newblock Cambridge Univ. Press, 2014.
\newblock [\epfmtdoi{10.1017/CBO9781139814782}]

\bibitem{reingold2013pseudorandomness}\bibhead{reingold2013pseudorandomness}
{\sc Omer Reingold, Thomas Steinke, and Salil Vadhan}: Pseudorandomness for
  regular branching programs via {F}ourier analysis.
\newblock In {\em Proc. Internat. Workshop on Approximation, Randomization, and
  Combinatorial Optimization (RANDOM'13)}, pp. 655--670. Springer, 2013.
\newblock [\epfmtdoi{10.1007/978-3-642-40328-6\_45}, \epfmt{arxiv}{1306.3004}]

\bibitem{simon1983tight}\bibhead{simon1983tight}
{\sc Hans-Ulrich Simon}: A tight {$\Omega(\log\log n)$-bound} on the time for
  parallel {RAM}'s to compute nondegenerated {B}oolean functions.
\newblock {\em Inf. Control}, 55(1--3):102--107, 1982.
\newblock See also
  \href{https://link.springer.com/chapter/10.1007/3-540-12689-9_124}{FCT'83}.
\newblock [\epfmtdoi{10.1016/S0019-9958(82)90477-6}]

\bibitem{tal2017tight}\bibhead{tal2017tight}
{\sc Avishay Tal}: Tight bounds on the {F}ourier spectrum of ${AC}^0$.
\newblock In {\em Proc. 32nd Computational Complexity Conf. (CCC'17)}, pp.
  15:1--15:31. Leibniz-Zentrum fuer Informatik, 2017.
\newblock [\epfmtdoi{10.4230/LIPIcs.CCC.2017.15}]

\bibitem{trevisan2013derandomized}\bibhead{trevisan2013derandomized}
{\sc Luca Trevisan and Tongke Xue}: A derandomized switching lemma and an
  improved derandomization of ${AC}^0$.
\newblock In {\em Proc. 28th IEEE Conf. on Computational Complexity (CCC'13)},
  pp. 242--247. IEEE Comp. Soc. Press, 2013.
\newblock [\epfmtdoi{10.1109/CCC.2013.32}]

\bibitem{vadhan2012pseudorandomness}\bibhead{vadhan2012pseudorandomness}
{\sc Salil Vadhan}: Pseudorandomness.
\newblock {\em Foundations and Trends in Theoret. Computer Sci.},
  7(1--3):1--336, Now Publ., 2012.
\newblock [\epfmtdoi{10.1561/0400000010}]

\bibitem{viola2009sum}\bibhead{viola2009sum}
{\sc Emanuele Viola}: The sum of ${D}$ small-bias generators fools polynomials
  of degree ${D}$.
\newblock {\em Comput. Complexity}, 18(2):209--217, 2009.
\newblock Preliminary version in
  \href{https://doi.org/10.1109/CCC.2008.16}{CCC'08}.
\newblock [\epfmtdoi{10.1007/s00037-009-0273-5}]

\end{thebibliography}

\end{document}
